% =====================================================================
% Va justo antes de \subsection{La ley de gravitación universal en
% términos del campo} (archivo gravitacion_campo.tex), dentro de
% \section{Ley de Gravitación Universal de Newton}.
% Requiere: amsmath, siunitx.
% =====================================================================

\subsection{La prueba de la Luna}
\label{subsec:prueba_luna}

Una bala disparada horizontalmente desde lo alto de una montaña cae hacia la
Tierra mientras avanza. Si su rapidez es suficiente, cae al mismo ritmo con
que la superficie se curva y nunca llega al suelo: queda en órbita. Newton se
preguntó si la Luna no sería precisamente eso, una bala de cañón muy lejana
que cae continuamente hacia la Tierra.

La Luna describe una órbita casi circular de radio
$r_L \approx \SI{3.84e8}{\meter}$ en un período $T \approx 27{,}3$~días
$\approx \SI{2.36e6}{\second}$. Su aceleración centrípeta, dirigida hacia el
centro de la Tierra, es
\begin{equation}
  a_L = \frac{v^{2}}{r_L} = \frac{4\pi^{2} r_L}{T^{2}}
  \approx \SI{2.72e-3}{\meter\per\second\squared} \,.
  \label{eq:aceleracion_luna}
\end{equation}
Comparada con la aceleración de caída libre en la superficie,
\begin{equation}
  \frac{g_0}{a_L} \approx \frac{\num{9.8}}{\num{2.72e-3}}
  \approx 3600 = 60^{2} \,,
  \label{eq:cociente_luna}
\end{equation}
y la Luna está a una distancia del centro de la Tierra de unos $60$ radios
terrestres, $r_L \approx 60\,R_T$. A una distancia $60$ veces mayor, la
aceleración es $60^{2}$ veces menor. Dicho de otra manera: en un segundo una
manzana cae $\SI{4.9}{\meter}$, mientras que la Luna se desvía de la línea
recta, hacia la Tierra, apenas $\SI{1.36}{\milli\meter}$.

Si suponemos que la aceleración de caída libre decae como $1/r^{n}$, la
predicción para la Luna es $g_0/60^{n}$. Con $n = 1$ se obtiene
$\SI{0.16}{\meter\per\second\squared}$, unas $60$ veces más de lo medido;
con $n = 3$, unas $60$ veces menos. Solo $n = 2$ concuerda con la
observación. La manzana y la Luna caen por la misma causa, debilitada con la
distancia según
\begin{equation}
  g(r) = g_0 \left( \frac{R_T}{r} \right)^{2} \,.
  \label{eq:g_inverso_cuadrado}
\end{equation}
En la siguiente subsección veremos por qué el exponente es precisamente
$2$.
